%2multibyte Version: 5.50.0.2960 CodePage: 1252

\documentclass{article}
\usepackage{amsfonts}

%%%%%%%%%%%%%%%%%%%%%%%%%%%%%%%%%%%%%%%%%%%%%%%%%%%%%%%%%%%%%%%%%%%%%%%%%%%%%%%%%%%%%%%%%%%%%%%%%%%%%%%%%%%%%%%%%%%%%%%%%%%%%%%%%%%%%%%%%%%%%%%%%%%%%%%%%%%%%%%%%%%%%%%%%%%%%%%%%%%%%%%%%%%%%%%%%%%%%%%%%%%%%%%%%%%%%%%%%%%%%%%%%%%
\usepackage{amsmath}

\setcounter{MaxMatrixCols}{10}
%TCIDATA{OutputFilter=LATEX.DLL}
%TCIDATA{Version=5.50.0.2960}
%TCIDATA{Codepage=1252}
%TCIDATA{<META NAME="SaveForMode" CONTENT="1">}
%TCIDATA{BibliographyScheme=Manual}
%TCIDATA{Created=Monday, September 22, 2003 09:40:11}
%TCIDATA{LastRevised=Tuesday, September 23, 2025 21:14:42}
%TCIDATA{<META NAME="GraphicsSave" CONTENT="32">}
%TCIDATA{<META NAME="DocumentShell" CONTENT="General\SW\Blank - Standard LaTeX Article">}
%TCIDATA{Language=American English}
%TCIDATA{CSTFile=LaTeX article (bright).cst}

\newtheorem{theorem}{Theorem}
\newtheorem{acknowledgement}[theorem]{Acknowledgement}
\newtheorem{algorithm}[theorem]{Algorithm}
\newtheorem{axiom}[theorem]{Axiom}
\newtheorem{case}[theorem]{Case}
\newtheorem{claim}[theorem]{Claim}
\newtheorem{conclusion}[theorem]{Conclusion}
\newtheorem{condition}[theorem]{Condition}
\newtheorem{conjecture}[theorem]{Conjecture}
\newtheorem{corollary}[theorem]{Corollary}
\newtheorem{criterion}[theorem]{Criterion}
\newtheorem{definition}[theorem]{Definition}
\newtheorem{example}[theorem]{Example}
\newtheorem{exercise}[theorem]{Exercise}
\newtheorem{lemma}[theorem]{Lemma}
\newtheorem{notation}[theorem]{Notation}
\newtheorem{problem}[theorem]{Problem}
\newtheorem{proposition}[theorem]{Proposition}
\newtheorem{remark}[theorem]{Remark}
\newtheorem{solution}[theorem]{Solution}
\newtheorem{summary}[theorem]{Summary}
\newenvironment{proof}[1][Proof]{\textbf{#1.} }{\ \rule{0.5em}{0.5em}}
\input{tcilatex}
\input{tcilatex}
\setlength{\textwidth}{6.5in}
\setlength{\textheight}{9in}
\setlength{\topmargin}{-.50in}
\setlength{\oddsidemargin}{.0in}
\setlength{\evensidemargin}{.0in}

\begin{document}


\begin{enumerate}
\item Consider an economy with a continuum of indentcal consumers and
normalize the mass of consumers to 1. Each agent is endowed with one unit of
time which is allocated across work and leisure. Preferences are $U\left(
c,\ell ,\bar{\ell}\right) =\ln c+\ln \ell +\ln \bar{\ell}$ where $c$ and $%
\ell $ are the individual's consumption and leisure and $\bar{\ell}$ is the
average leisure across the population. Note that because any individual is
very small relative to the population, each consumer will treat $\bar{\ell}$
as given. There is an externality in that the utility of an individual is
influenced by others' leisure. A representative firm behaves competitively
and hires labor, $n$, to maximize profits with production technology given
by $y=n^{\alpha }$.

\begin{enumerate}
\item Determine the Pareto optimal leisure.

\item Determine leisure in a competitive equilibrium.

\item Now suppose that the government taxes labor income at rate $\tau $ and
returns the revenue lump sum to the representative agents as a subsidy, $%
\varsigma $, subject to a balanced budget constraint. Write down the agent's
optimization problem and find the first order condition (or conditions
depending on how you have written the problem). Find and impose any
equilibrium conditions that are required so that you can solve for $\ell $
in a competitive equilibrium. Solve for $\ell .$

\item Consider the same situation as in part (c) except there are now $k\geq
2$ agents so that $\bar{\ell}=\frac{1}{k}\underset{i=1}{\overset{k}{\dsum }}%
\ell _{1}.$  Here $\ell _{1}$ is the amount of leisuer for the agent we are
considering and  the other $\ell _{i}$ are  leisure for the other agents.
Note that $k$ may now be small enough to be influenced by the agent's
decisions. All $\ell _{i}$ for $i>1$ are taken as given by the agent (agent
1) but all agents are agaid indentcal. Find $\ell $ in equibrium and discuss
how the relationship between the social planner's allocation and the
competetive equilbrium changes as  $k$ grows. Give the intuition. 
\end{enumerate}

\item ..., Solution is: $\left\{ 
\begin{array}{ccc}
\emptyset  & \text{if} & \tau =1\wedge n=-1\vee -2n+n\tau -1=0\wedge n\neq
-1\vee \tau \neq 1\wedge n=-1\vee \left( \frac{1}{2n-n\tau +1}+\frac{n}{%
2n-n\tau +1}=1\vee \frac{n+1}{2n-n\tau +1}=0\right) \wedge 2n-n\tau +1\neq
0\wedge n\neq -1 \\ 
\left\{ {}\right\}  & \text{if} & -2n+n\tau -1\neq 0\wedge \frac{n+1}{%
2n-n\tau +1}\in 
%TCIMACRO{\U{2102} }%
%BeginExpansion
\mathbb{C}
%EndExpansion
\setminus \left\{ 0,1\right\} \wedge n\neq -1%
\end{array}%
\right. \allowbreak $%
\begin{equation*}
\frac{1-\tau }{1-\ell }=\frac{1}{\ell }
\end{equation*}%
\begin{equation*}
\ell =\frac{1}{2-\tau }
\end{equation*}

\begin{enumerate}
\item Set $\tau $ to make these equal. 
\begin{equation*}
\frac{2}{\alpha +2}=\frac{1}{1-\tau +1}
\end{equation*}%
\begin{equation*}
\frac{2}{\alpha +2}=\frac{1}{2-\tau }
\end{equation*}%
\begin{eqnarray*}
\tau  &=&1-\frac{1}{2}\alpha  \\
&&
\end{eqnarray*}
: $1-\frac{1}{2}\alpha $
\end{enumerate}

\item .Now look at competitive equilibrium. \ Agent's problem is 
\begin{equation*}
U\left( c,\ell ,\bar{c}\right) =\ln c+\ln \ell +\ln \left( \frac{\ell
+\left( n-1\right) \ell _{o}}{n}\right) 
\end{equation*}%
st. $c=\omega n$%
\begin{equation*}
U\left( c,\ell ,\bar{c}\right) =\ln c+\ln \ell +\ln \left( \frac{\ell
+\left( n-1\right) \ell _{o}}{n}\right) 
\end{equation*}%
\begin{equation*}
U\left( c,\ell ,\bar{c}\right) =\ln \left( w\left( 1-\ell \right) \right)
+\ln \ell +\ln \left( \frac{\ell +\ell _{o}}{2}\right) 
\end{equation*}%
\begin{equation*}
U\left( c,\ell ,\bar{c}\right) =\ln \left( w\left( 1-\ell \right) \right)
+\ln \ell +\ln \left( \frac{\ell +\ell _{o}}{2}\right) 
\end{equation*}%
\begin{equation*}
\frac{1}{1-\ell }=\frac{1}{\ell }+\frac{1}{\ell +\left( n-1\right) \ell _{o}}
\end{equation*}%
\begin{equation*}
\frac{1}{1-\ell }=\frac{1}{\ell }+\frac{1}{\ell +\left( n-1\right) \ell }
\end{equation*}%
\begin{equation*}
\frac{1}{1-\ell }=\frac{1}{\ell }+\frac{1}{n\ell }
\end{equation*}%
, Solution is: $\left\{ 
\begin{array}{ccc}
\emptyset  & \text{if} & n=0\vee n=-\frac{1}{2} \\ 
\left\{ {}\right\}  & \text{if} & \frac{1}{2n+1}\left( n+1\right) \in 
%TCIMACRO{\U{2102} }%
%BeginExpansion
\mathbb{C}
%EndExpansion
\setminus \left\{ 0,1\right\} \wedge n\neq 0\wedge n\neq -\frac{1}{2}%
\end{array}%
\right. \allowbreak \frac{1}{2n+1}\left( n+1\right) $%
\begin{equation*}
\frac{1}{1-\ell }=\frac{1}{\ell }+\frac{1}{2\ell }
\end{equation*}%
\begin{equation*}
\frac{2}{1-\ell }=\frac{2}{\ell }+\frac{1}{\ell }
\end{equation*}%
\begin{equation*}
\ell 2=3\left( 1-\ell \right) 
\end{equation*}%
\begin{equation*}
\ell =\frac{3}{5}
\end{equation*}%
\begin{equation*}
U\left( c,\ell ,\bar{c}\right) =\ln \left( w\left( 1-\ell \right) \right)
+\ln \ell +\ln \left( \frac{\ell +\ell _{a}+\ell _{b}}{3}\right) 
\end{equation*}
\end{enumerate}

\end{document}
